% !TeX program = lualatex
% Compile twice: lualatex open_problems_500.tex
% All references and prose are embedded; no BibTeX database or external inputs.
\documentclass[11pt,a4paper]{article}
\usepackage[margin=24mm,headheight=14pt]{geometry}
\usepackage{fontspec}
\setmainfont{Latin Modern Roman}
\setsansfont{Latin Modern Sans}
\setmonofont{Latin Modern Mono}
\usepackage{amsmath,amssymb,mathtools}
\usepackage{unicode-math}
\setmathfont{Latin Modern Math}
\usepackage{microtype}
\usepackage{enumitem,xurl,needspace}
\usepackage[dvipsnames]{xcolor}
\usepackage{hyperref}
\hypersetup{unicode=true,colorlinks=true,linkcolor=MidnightBlue,urlcolor=MidnightBlue,
 pdftitle={500 Mathematical Problem Statements},pdfauthor={Source-annotated compilation},bookmarksnumbered=true}
\usepackage{bookmark}
\usepackage{tocloft}
\setlength{\cftsecnumwidth}{3em}
\cftsetpnumwidth{2.5em}
\cftsetrmarg{4em}
\usepackage{titlesec}
\titleformat{\section}{\Large\bfseries\raggedright}{\thesection}{0.7em}{}
\usepackage{fancyhdr}
\pagestyle{fancy}\fancyhf{}
\fancyhead[L]{\small\sffamily Mathematical problem statements}
\fancyhead[R]{\small Source edition 22}
\fancyfoot[C]{\thepage}
\setlength{\parindent}{0pt}
\setlength{\parskip}{5pt plus 1pt}
\setlength{\emergencystretch}{3em}
\setcounter{tocdepth}{1}
\setcounter{secnumdepth}{1}
\setlist[itemize]{leftmargin=3em,itemsep=5pt,topsep=4pt,parsep=0pt}
\allowdisplaybreaks
\newcommand{\N}{\mathbb{N}}
\newcommand{\Z}{\mathbb{Z}}
\newcommand{\Q}{\mathbb{Q}}
\newcommand{\R}{\mathbb{R}}
\newcommand{\C}{\mathbb{C}}
\newcommand{\F}{\mathbb{F}}
\newcommand{\A}{\mathbb{A}}
\newcommand{\PP}{\mathbb P}
\newcommand{\E}{\mathbb{E}}
\newcommand{\eps}{\varepsilon}
\newcommand{\dd}{\,\mathrm d}
\newcommand{\sref}[2]{\hyperlink{source:#1:#2}{[S#2]}}
\newcommand{\eref}[2]{\hyperlink{extra:#1:#2}{[E#2]}}
% Turn the next switch off for a shorter reading copy. The quotations preserve
% every exactTarget field, including qualifications omitted from a short summary.
\newif\ifshowcatalogue
\showcataloguetrue
\newcommand{\cataloguescope}[1]{\ifshowcatalogue
 \paragraph{Exact scope in the supplied catalogue.}
 \begin{quote}\small #1\end{quote}\fi}
% Standalone notebook for original catalogue section 162.
% Compile twice with: lualatex 162.tex
\numberwithin{equation}{section}
\hypersetup{pdftitle={Research attempt -- problem 162},pdfauthor={Research notebook; original compilation retained}}
\fancyhead[L]{\small\sffamily Research notebook}
\fancyhead[R]{\small\sffamily Problem 162 | 27 September 2026}
\begin{document}
\setcounter{section}{161}
% ================================================================
% problemId: problem.keating-snaith-conjecture-on-moments-of-the-riemann-zeta-function
% source releaseRank: 162; source releaseStatus: open_with_solved_subcases
% exactTarget SHA-256: a58847d6e4e8f4d9e6ae63abf008de702afe1e1b1b5ca44de173c99426bf1379
\clearpage
\section{\texorpdfstring{Keating-Snaith Conjecture on Moments of the Riemann Zeta Function}{Keating-Snaith Conjecture on Moments of the Riemann Zeta Function}}\label{162}
\subsection{Definitions and mathematical statement}
For fixed real $k>0$, define $I_k(T)=\int_0^T|\zeta(\tfrac12+it)|^{2k}\dd t$. Let $G$ be the standard Barnes $G$-function, normalized by $G(1)=1$ and its usual canonical product, with $G(z+1)=\Gamma(z)G(z)$. Define
\[
 a_k=\prod_{p\ \mathrm{prime}}(1-p^{-1})^{k^2}
 \sum_{m=0}^{\infty}\left(\frac{\Gamma(m+k)}{\Gamma(k)m!}\right)^2p^{-m}.
\]
The assertion is
$I_k(T)\sim a_kG(1+k)^2G(1+2k)^{-1}T(\log T)^{k^2}$ as $T\to\infty$. The parameter $k$ is fixed before this limit, and the target includes the leading constant, not just the logarithmic exponent. The functional equation alone is not being used as a unique definition of $G$.
\par\smallskip\textit{Formulation and concept references:} \sref{162}{1}.
\cataloguescope{For every fixed real k>0, let I\_k(T)=integral\_0\textasciicircum{}T |zeta(1/2+it)|\textasciicircum{}(2k) dt. Prove that I\_k(T) is asymptotic to a\_k G(1+k)\textasciicircum{}2/G(1+2k) times T(log T)\textasciicircum{}(k\textasciicircum{}2), where G is the Barnes G-function and a\_k=product\_p (1-1/p)\textasciicircum{}(k\textasciicircum{}2) sum\_(m>=0) (Gamma(m+k)/(Gamma(k)m!))\textasciicircum{}2 p\textasciicircum{}(-m).}
\subsection{Short English statement}
Do all positive real moments of the zeta function on its central line have the precise leading asymptotic predicted by the arithmetic product and random-matrix factor?
% BEGIN RESEARCH ADDITION ATTEMPT-162
\subsection{Research attempt}

\paragraph{Current frontier.}
The exact leading asymptotic is classical for $k=1,2$, whereas the
all-positive-real-$k$ assertion remains the conjecture stated in
\sref{162}{1}. Florea's 2026 account distinguishes the leading-constant
problem from sharp order-of-magnitude bounds \eref{162}{1}, Sections
1--3. Harper proves the upper order $T(\log T)^{k^2}$ for each fixed
$k\ge0$ assuming RH, without identifying the leading constant
\eref{162}{2}. The hybrid Euler--Hadamard approach separates prime and
zero factors but its requisite moment splitting is conjectural in general
\eref{162}{3}. These distinctions are retained below; RH and a
random-matrix model are not assumed to be theorems.

\paragraph{Attack route.}
Direct mean-square evaluation of a polynomial approximating $\zeta^k$
loses control once its length becomes too large. A hybrid factorization
could preserve the arithmetic Euler factor while transferring only the
remaining component to a zero statistic. The route pursued here proves
the arithmetic-side statements with a slowly growing cutoff, then writes
the remaining requirement as a \emph{weighted} moment. The important
question is whether the prime weight can legitimately be decoupled from
that residual; its variance will show why a naive Cauchy--Schwarz step
is inadequate.

\paragraph{Attempt: checking the arithmetic constant.}
Put $d_k(p^m)=(k)_m/m!$, for fixed $k>0$. On any compact interval of
such $k$, its polynomial growth in $m$ implies
\[
 \sum_{m\ge0}d_k(p^m)^2p^{-m}=1+\frac{k^2}{p}+O_k(p^{-2}).
\]
Multiplication by $(1-p^{-1})^{k^2}$ cancels the term of order $p^{-1}$.
The resulting local factors are positive and $1+O_k(p^{-2})$, so the
Euler product $a_k$ converges to a positive number. This argument also
justifies finite-prime product manipulations made below; it does not
assert convergence of the Euler product for $\zeta$ on the critical line.

For $k=1$, $a_1=1$ and $G(2)^2/G(3)=1$. For $k=2$,
\[
 \sum_{m\ge0}(m+1)^2x^m=\frac{1+x}{(1-x)^3},\qquad
 a_2=\prod_p(1-p^{-2})=\zeta(2)^{-1},
 \qquad\frac{G(3)^2}{G(5)}=\frac1{12}.
\]
Thus the proposed fourth-moment coefficient is $1/(2\pi^2)$, as required
by the established case. The next integer Barnes factor is
$G(4)^2/G(7)=1/8640$. These identities are independently checked by the
script; their correctness supplies no estimate for the sixth moment.

\paragraph{Attempt: an explicit mean-square barrier.}
For $A(t)=\sum_{n\le N}a_n n^{-it}$ and
$S=\sum_{n\le N}|a_n|^2$, direct integration gives
\begin{equation}\label{162:mean-square}
 \left|\frac1T\int_T^{2T}|A(t)|^2\dd t-S\right|
 \le\frac{4N(1+\log N)}{T}S.
\end{equation}
To prove it, the off-diagonal integral has absolute value at most
$2/(T|\log(n/m)|)$. For $n>m$,
$\log(n/m)\ge(n-m)/n\ge(n-m)/N$. Hence the total error is at most
\[
 \frac{4N}{T}\sum_{d=1}^{N-1}\frac1d
       \sum_{m\le N-d}|a_ma_{m+d}|
 \le \frac{4N(1+\log N)}{T}S,
\]
where the inner sum is bounded by $S$ using Cauchy--Schwarz. All sums
are finite. Thus $N\log N=o(T)$ is a sufficient diagonal regime for this
particular estimate. It is not a necessary condition for all possible
mean-value methods. In particular, the inequality does not dismiss the
more delicate off-diagonal methods that establish the fourth moment.
It merely shows exactly why extending this direct argument to long
polynomials would be unjustified.

\paragraph{Attempt: a rigorous finite-prime factor.}
For $y\ge2$ define the finite, never-vanishing product
\[
 P_y(t)=\prod_{p\le y}(1-p^{-1/2-it})^{-1},\qquad
 \langle F\rangle_T=\frac1T\int_T^{2T}F(t)\dd t.
\]
For each fixed $y$ and $k>0$,
\begin{equation}\label{162:fixed-y}
 \lim_{T\to\infty}\langle|P_y|^{2k}\rangle_T
 =A_y(k):=\prod_{p\le y}\sum_{m\ge0}d_k(p^m)^2p^{-m}.
\end{equation}
Here is a justification not invoking a random-prime assumption. A nonzero
integer combination of the numbers $\log p$, $p\le y$, cannot be zero,
by unique factorization. The average of each nonconstant character of
the resulting torus flow therefore tends to zero by direct integration.
Approximation by trigonometric polynomials extends this to continuous
functions on the finite torus. The integrand in question is continuous
because $p^{-1/2}<1$. On each circle the absolutely convergent binomial
series for $(1-z)^{-k}$ and orthogonality give the local sum in
\eqref{162:fixed-y}. The finite torus integral factors.

The convergent Euler product and Mertens' product theorem give
\begin{equation}\label{162:Ay}
 A_y(k)\sim a_k(e^\gamma\log y)^{k^2}
 \qquad(y\to\infty).
\end{equation}
These familiar prime-factor asymptotics are also the arithmetic input in
\eref{162}{3}; the product $P_y$ used here is an elementary finite product,
not an assertion of the full smoothed hybrid formula from that paper.

A diagonal choice now provides $y=y(T)\to\infty$ with
$\log y=o(\log T)$ for which
\eqref{162:fixed-y} holds with relative error $o(1)$ simultaneously for
$k$ and $2k$. Explicitly, choose a sequence $y_j\to\infty$, then a
strictly increasing sequence $T_j$ so that both fixed-$y_j$ relative
errors are at most $1/j$ for every $T\ge T_j$, and also
$\log y_j\le\sqrt{\log T_j}$. Put $y(T)=y_j$ between successive thresholds.
This is an existence construction, not a quantitative cutoff such as
$y=T^\theta$.

\paragraph{Attempt: the remaining weighted statement.}
For this cutoff write $b_T=e^\gamma\log y(T)$ and set
\[
 Q_T(t)=\frac{|P_{y(T)}(t)|^{2k}}
                   {\langle|P_{y(T)}|^{2k}\rangle_T},\qquad
 D_T(t)=\left(\frac{b_T}{\log T}\right)^{k^2}
          \left|\frac{\zeta(1/2+it)}{P_{y(T)}(t)}\right|^{2k}.
\]
These are nonnegative measurable functions and $\langle Q_T\rangle_T=1$.
There is an exact identity
\begin{equation}\label{162:weighted-identity}
 \frac{\langle|\zeta(1/2+it)|^{2k}\rangle_T}
        {a_k(\log T)^{k^2}}
 =\frac{\langle|P_{y(T)}|^{2k}\rangle_T}{a_k b_T^{k^2}}
       \langle Q_TD_T\rangle_T.
\end{equation}
The prefactor tends to one. Thus the local moment asymptotic reduces to
\begin{equation}\label{162:missing-weighted}
 \langle Q_TD_T\rangle_T\longrightarrow
 g_k:=\frac{G(1+k)^2}{G(1+2k)}.
\end{equation}
For example, the two genuinely additional statements
$\langle D_T\rangle_T\to g_k$ and
$\langle(Q_T-1)D_T\rangle_T\to0$ would suffice. Neither has been proved
in this argument. Summing local $[T,2T]$ asymptotics over dyadic intervals
would then give the requested integral from $0$ to $T$: finitely many
initial intervals contribute $o(T(\log T)^{k^2})$, and the geometric tail
of the interval lengths controls the remaining error. Conversely the
printed global asymptotic implies the local one by subtraction.

\paragraph{Stress test: the weight is not harmless.}
The same fixed-prime computation yields
\begin{equation}\label{162:variance}
 \langle Q_T^2\rangle_T
 \sim\frac{a_{2k}}{a_k^2}b_T^{2k^2}\longrightarrow\infty.
\end{equation}
Consequently $\langle D_T\rangle_T\to g_k$ alone says nothing about the
covariance in \eqref{162:missing-weighted}. A Cauchy--Schwarz argument
would need, for example,
$\|D_T-g_k\|_2=o(b_T^{-k^2})$, an extremely strong concentration
statement not supplied by a moment model. Nonnegative random variables
with the right separate means can still have substantial positive
correlation; the exact identity has not removed this arithmetic obstacle.

Weak distributional convergence is likewise insufficient. To see this
without making a claim about $\zeta$, let $L\to\infty$ and start with the
model $Z_L=\exp(\sqrt{L/2}\,N)$ for a standard Gaussian $N$.
Its $2k$-th moment is $e^{k^2L}$. With probability
$e^{-(2kb-k^2)L}$, for any $b>k/2$, replace it by $e^{bL}$.
This probability tends to zero, so the normalized logarithm has the same
weak Gaussian limit, but the added contribution to its $2k$-th moment
is exactly $e^{k^2L}$. The leading coefficient doubles asymptotically.
This construction falsifies the proposed inference from a central-limit
law to the desired moment constant.

\paragraph{Outcome.}
\textbf{Outcome: Exploratory approach with unresolved gap.}

The finite-product moment evaluation, diagonal cutoff construction,
weighted identity, and variance calculation are justified. They isolate
a weighted prime--residual correlation requirement and rule out two
insufficient transfers: a bare Gaussian limit and an unweighted residual
mean. The missing statement \eqref{162:missing-weighted} still contains
the difficult zero statistics and their correlation with primes. No
leading constant beyond the established cases has been proved here.
% END RESEARCH ADDITION ATTEMPT-162
\subsection{Sources}
\begin{itemize}\raggedright
\item[\textnormal{[S1]}]\hypertarget{source:162:1}{} J. P. Keating and N. C. Snaith, "Random Matrix Theory and zeta(1/2+it)," Communications in Mathematical Physics 214 (2000), 57-89.
\url{https://empslocal.ex.ac.uk/people/staff/mrwatkin/zeta/keating-snaith.pdf}.
{\small\textit{Catalogue formulation source.}}
% BEGIN RESEARCH ADDITION SOURCES-162
\item[\textnormal{[E1]}]\hypertarget{extra:162:1}{} Alexandra Florea, ``A survey of moment bounds for $\zeta(s)$: From Heath-Brown's work to the present,'' \emph{J. London Math. Soc.} 113 (2026), no. 1, e70376. Sections 1--3.
\url{https://doi.org/10.1112/jlms.70376}.
{\small\textit{Current distinction between exact moment asymptotics and moment-size bounds.}}
\item[\textnormal{[E2]}]\hypertarget{extra:162:2}{} Adam J. Harper, ``Sharp conditional bounds for moments of the Riemann zeta function,'' arXiv:1305.4618 (2013). Main theorem, under the Riemann Hypothesis.
\url{https://arxiv.org/abs/1305.4618}.
{\small\textit{Conditional upper bounds, with no asserted leading constant.}}
\item[\textnormal{[E3]}]\hypertarget{extra:162:3}{} S. M. Gonek, C. P. Hughes and J. P. Keating, ``A hybrid Euler--Hadamard product for the Riemann zeta function,'' \emph{Duke Math. J.} 136 (2007), 507--549; arXiv:math/0511182. Introduction and prime/zero moment-splitting discussion.
\url{https://www.sas.rochester.edu/mth/people/faculty/gonek-steve/assets/pdf/hybridformula.pdf}.
{\small\textit{Motivation for separating arithmetic and zero factors; no conjectural splitting is used as a theorem.}}
% END RESEARCH ADDITION SOURCES-162
\end{itemize}

\end{document}
